\documentclass[11pt,a4paper]{article}

% ---------------------------------------------------------
% PREAMBLE (stable, warning-free, Overleaf-safe)
% ---------------------------------------------------------
\usepackage[utf8]{inputenc}
\usepackage[T1]{fontenc}
\usepackage{lmodern}
\usepackage{amsmath, amssymb, amsfonts}
\usepackage{bm}
\usepackage{geometry}
\usepackage{graphicx}
\usepackage{hyperref}
\usepackage{cite}
\usepackage{authblk}
\usepackage{physics}
\usepackage{mathtools}

\geometry{margin=2.4cm}

\title{\textbf{Companion XL: CMB Polarisation in the Relaxation-Driven Cyclic Cosmology}}
\author{Michael Lehmann \\ \texttt{mi.lehmann@gmx.de}}
\date{April 2026}

\begin{document}
\maketitle

\begin{abstract}
The polarisation of the cosmic microwave background (CMB) encodes the temporal 
structure of the gravitational field at the epoch of recombination and during the 
subsequent evolution of the Universe. In the standard cosmological model, E-mode 
polarisation arises from scalar perturbations, while B-modes are generated by 
tensor perturbations or by lensing of E-modes. In the Relaxation-Driven Cyclic 
Cosmology (RDCC), the scale-dependent evolution of the gravitational potential 
modifies the quadrupole anisotropy that sources polarisation, leading to 
characteristic signatures in both E- and B-modes. This Companion develops a 
continuous description of CMB polarisation in RDCC, deriving how the relaxation 
scale influences the generation of E-modes, the leakage into B-modes, and the 
late-time evolution of polarisation through the integrated Sachs–Wolfe effect. 
The resulting framework provides a new observational window into the temporal 
architecture of the RDCC gravitational field.
\end{abstract}
% ---------------------------------------------------------
\section{Introduction}

CMB polarisation arises from Thomson scattering of photons off free electrons in 
the presence of a quadrupole temperature anisotropy. The structure of this 
quadrupole is determined by the evolution of the gravitational potential and the 
velocity field of the photon–baryon fluid. In the standard cosmological model, the 
potential evolves smoothly, and the resulting polarisation pattern is dominated by 
E-modes, with B-modes generated only by tensor perturbations or by lensing.

In RDCC, the gravitational potential evolves differently. The relaxation mechanism 
suppresses the decay of small-scale modes and enhances the coherence of large-scale 
modes. This modifies the quadrupole anisotropy at recombination and during the 
subsequent evolution of the Universe. As a result, the polarisation pattern 
acquires a distinctive RDCC signature that reflects the temporal structure of the 
gravitational field.
% ---------------------------------------------------------
\section{The RDCC Gravitational Potential and the Quadrupole Source}

The generation of CMB polarisation is governed by the quadrupole moment of the 
photon distribution. This quadrupole is sourced by the time evolution of the 
gravitational potential. In RDCC, the potential takes the form
\begin{equation}
    \Phi_{\mathrm{RDCC}}(k,a)
    = \Phi(k,a)
      \left[
        1 - \chi_{\Phi}
        \left(\frac{k}{k_{\mathrm{rel}}}\right)^{\alpha}
      \right].
\end{equation}

The quadrupole source term becomes
\begin{equation}
    \Theta_{2,\mathrm{RDCC}}(k,a)
    = \Theta_{2}(k,a)
      - \chi_{2}
        \left(\frac{k}{k_{\mathrm{rel}}}\right)^{\alpha}
        \frac{\partial \Phi}{\partial a}.
\end{equation}

Small-scale modes contribute less to the quadrupole, while large-scale modes 
retain their coherence. This imbalance reshapes the polarisation pattern at 
recombination.
% ---------------------------------------------------------
\section{E-Mode Polarisation in RDCC}

E-mode polarisation arises from the curl-free component of the polarisation field. 
In RDCC, the suppression of small-scale potential evolution reduces the amplitude 
of high-$\ell$ E-modes. The E-mode power spectrum becomes
\begin{equation}
    C_{\ell}^{EE,\mathrm{RDCC}}
    = C_{\ell}^{EE}
      \left[
        1 - \chi_{E}
        \left(\frac{\ell}{\ell_{\mathrm{rel}}}\right)^{\alpha}
      \right],
\end{equation}
where $\ell_{\mathrm{rel}}$ corresponds to the angular scale of the relaxation 
scale $k_{\mathrm{rel}}$.

This modification produces a smoother E-mode spectrum with reduced small-scale 
power and enhanced large-scale coherence.
% ---------------------------------------------------------
\section{B-Mode Polarisation and RDCC-Induced Leakage}

B-mode polarisation is generated by tensor perturbations or by lensing of E-modes. 
In RDCC, the altered evolution of the gravitational potential introduces an 
additional source of B-modes. The leakage from E- to B-modes becomes
\begin{equation}
    C_{\ell}^{EB,\mathrm{RDCC}}
    \propto
    \left(
        \frac{\partial \Phi_{\mathrm{RDCC}}}{\partial a}
    \right)
    C_{\ell}^{EE}.
\end{equation}

Because the relaxation mechanism modifies the time derivative of the potential, 
this leakage acquires a scale dependence that reflects the relaxation scale. The 
resulting B-mode spectrum is not dominated by primordial tensor modes but by the 
temporal structure of the RDCC gravitational field.

This provides a new observational signature that distinguishes RDCC from the 
standard cosmological model.
% ---------------------------------------------------------
\section{Late-Time Evolution and the Integrated Sachs--Wolfe Contribution}

The integrated Sachs--Wolfe (ISW) effect contributes to CMB polarisation through 
the late-time evolution of the gravitational potential. In RDCC, the potential 
evolves differently at late times, leading to a modified ISW contribution to both 
E- and B-modes.

The ISW-induced polarisation becomes
\begin{equation}
    \Delta P_{\mathrm{ISW,RDCC}}
    \propto
    \int da\,
    \frac{\partial \Phi_{\mathrm{RDCC}}}{\partial a}
    \Theta_{1}(a),
\end{equation}
where $\Theta_{1}$ is the dipole moment of the photon distribution.

Because the relaxation mechanism suppresses the decay of small-scale modes, the 
ISW contribution becomes smoother and more coherent, enhancing large-scale 
polarisation.
% ---------------------------------------------------------
\section{Conclusion}

CMB polarisation in the Relaxation-Driven Cyclic Cosmology reflects the 
scale-dependent evolution of the gravitational potential. The relaxation mechanism 
modifies the quadrupole anisotropy at recombination, the generation of E-modes, 
the leakage into B-modes, and the late-time ISW contribution. These effects 
produce distinctive signatures in the polarisation power spectra that provide a 
direct observational probe of the relaxation scale and the temporal structure of 
the gravitational field. The present Companion establishes the theoretical 
foundation for future studies of CMB secondary anisotropies, lensing-induced 
polarisation, and the interplay between scalar and tensor modes in RDCC.

% ========================
% References
% ========================

\begin{thebibliography}{10}

\bibitem{Flagship} Lehmann, M. (2026). Relaxation-Driven Cyclic Cosmology (RDCC): A Minimal CPT-Symmetric Framework with a Single Parameter. Flagship Paper. Zenodo: \url{https://zenodo.org/records/19744958} (v43.0 or later).

\bibitem{Zenodo} The complete RDCC companion series (I--LVII) is available at the same Zenodo record above.

% Thematisch relevante Companions
\bibitem{CompXVI} Companion XVI: Boltzmann Code and Modified Hierarchy.
\bibitem{CompXXXI} Companion XXXI: RDCC Halo Model and Non-Linear Structure.
\bibitem{CompXXXII} Companion XXXII--XXXIX: Galaxy--Halo Connection, Cosmic Web, Lensing, RSD, ISW, tSZ, Rotation Curves, CGM.

% Wichtige externe Referenzen
\bibitem{Planck2020} Planck Collaboration (2020), Astron. Astrophys. 641, A6.
\bibitem{DESI2024} DESI Collaboration (2024/2025), arXiv: [aktuelle $f\sigma_8$/BAO-Paper].
\bibitem{JWST2024} Maiolino et al. (2024), Nature (JWST high-$z$ galaxies/quasars).
\bibitem{Kaiser1987} Kaiser, N. (1987), Mon. Not. R. Astron. Soc. 227, 1 (RSD).
\bibitem{Bartelmann2001} Bartelmann, M. \& Schneider, P. (2001), Phys. Rep. 340, 291 (Weak Lensing Review).
\bibitem{HuOkamoto2002} Hu, W. \& Okamoto, T. (2002), Astrophys. J. 574, 566 (CMB Lensing).
\bibitem{LewisChallinor2006} Lewis, A. \& Challinor, A. (2006), Phys. Rep. 429, 1 (CMB Lensing Review).
\bibitem{NFW1997} Navarro, J. F., Frenk, C. S. \& White, S. D. M. (1997), Astrophys. J. 490, 493 (Halo Profiles).

\end{thebibliography}

\end{document}
